\begin{table}[t]
  \centering
  \caption{\textbf{視点分布を切り分けるための比較計画。}
  本稿の時点ですべて未実施である。
  各条件は同じ再構成シーン、同じ教師 schema、同じモデル、
  同じ初期重み、同じ更新回数を持ち、
  動かす変数を一行に一つだけにする。
  }
  \label{tab:ablation_plan}
  \small
  \begin{tabularx}{\linewidth}{@{}p{0.05\linewidth}p{0.28\linewidth}X@{}}
    \toprule
    & 学習データ & 切り分けるもの \\
    \midrule
    A & 実写のみ & 実写学習の基準 \\
    B & 合成・視点狭域のみ（同枚数） & 制限された視点での基準 \\
    C & 合成・視点広域のみ（B と同枚数） & 視点範囲だけの効果 \\
    D & 同じ実写量 + 合成・狭域（同じ合成枚数） & 実写を使う条件の基準 \\
    E & 同じ実写量 + 合成・広域（D と同合成枚数） & 実写下での視点範囲の効果 \\
    \bottomrule
  \end{tabularx}
\end{table}
